
\subsubsection{Interpretation}

The original hypothesis predicted that NFN would require 50\% fewer samples than MLP to achieve equivalent performance. The experimental results reveal a stronger finding: the comparison is not quantitative but categorical. The MLP baseline does not merely require more data---it fails completely at all sample sizes tested.

The MLP faces ~270K input features with only 500--5000 training samples---a severely underdetermined problem. Each hidden layer has N! equivalent permutations. Without equivariance, the MLP cannot generalize across permutation-equivalent inputs.

The NFN succeeds by constraining the hypothesis space through weight sharing and architectural invariance, reducing effective dimensionality. The statistics baseline succeeds via aggressive projection (270K $\rightarrow$ 63 features), trading expressiveness for guaranteed invariance.

\subsubsection{Theoretical Interpretation}

The MLP failure arises from a combination of two factors:
\begin{enumerate}
\item \textbf{Dimensionality}: Insufficient samples for the input dimensionality
\item \textbf{Symmetry}: Cannot generalize across permutation-equivalent configurations
\end{enumerate}

The equivariant NFN addresses both by encoding permutation structure architecturally, enabling learning with limited data.

\subsubsection{Connections to Prior Work}

The statistics baseline result ($R^2$ = 0.9995) confirms and extends Unterthiner et al. (2020), who achieved $R^2$ > 0.98. The NFN results align with Zhou et al. (2023) in demonstrating effective equivariant weight processing. The categorical failure of the MLP provides empirical support for the permutation symmetry analysis in Ainsworth et al. (2023).

\subsubsection{Limitations}

\textbf{L1: Synthetic Model Zoo}. Experiments used synthetically generated models rather than the full Zenodo Model Zoo (157GB). Accuracy distributions may be more uniform than real-world zoos; the statistics baseline may overperform relative to expectations on real data. Results are directionally valid but require replication on real model collections.

\textbf{L2: Single Architecture}. All experiments used ResNet-20. Cross-architecture generalization (ResNet-50, ViT, mixed-architecture zoos) was not tested. NFN supports arbitrary architectures by design; ResNet-20 is representative but not exhaustive.

\textbf{L3: MLP Configuration}. The MLP used a standard 2-layer, 256-unit architecture. Alternative configurations (deeper networks, stronger regularization, dimensionality reduction preprocessing) were not exhaustively tested. Given the categorical nature of the failure ($R^2$ << 0), incremental improvements are unlikely to close the 2.5 $R^2$ gap.

\textbf{L4: Preliminary N=5000 Results}. The H-C1 convergence test used only seed 0. Multi-seed validation would strengthen confidence in the convergence failure, though the directional result is clear.

\subsubsection{Implications}

Weight-space learning systems must either use permutation-invariant features or equivariant architectures. Treating weights as arbitrary vectors fails regardless of dataset size within practical ranges.
